\subsubsection{Closest pair of points}

\paragraph{Idea intuitiva.}
Encuentra el par de puntos mas cercano en $O(n \log n)$ con \textbf{divide y venceras}. Algoritmo:
\begin{itemize}\setlength\itemsep{0pt}
	\item Sort por $x$.
	\item Dividir en mitades; recursivo.
	\item Combinar: dado el minimo $d$ de las dos mitades, strip = puntos en $[midX - d, midX + d]$; ordenar por $y$; cada punto solo se compara con los siguientes hasta $y$ difiera en $< d$ (maximo 7).
\end{itemize}

\paragraph{Propiedades clave (invariantes).}
\begin{itemize}\setlength\itemsep{0pt}
	\item Requiere ordenar por $y$ despues (en el snippet se hace durante el merge).
	\item Strip tiene a lo sumo $O(n)$ puntos en total, pero la comparacion es $O(1)$ amortizada por punto.
	\item Resultado en \texttt{double}.
\end{itemize}

\paragraph{Complejidad.}
\begin{itemize}\setlength\itemsep{0pt}
	\item $O(n \log n)$.
\end{itemize}

\paragraph{Donde tocar para modificar.}
\begin{itemize}\setlength\itemsep{0pt}
	\item \textbf{Reconstruir el par}: aniadir indices en la recursion.
	\item \textbf{Distancia maxima}: cambiar \texttt{min} por \texttt{max} (mas facil; no requiere D\&C).
	\item \textbf{Puntos duplicados}: aniadir check especial (distancia 0).
\end{itemize}

\paragraph{Trampas y errores frecuentes.}
\begin{itemize}\setlength\itemsep{0pt}
	\item El merge requiere ordenar por $y$; el snippet lo hace in-place con un \texttt{tmp}.
	\item Si se hace en cada recursion, es $O(n \log n)$ extra.
\end{itemize}

\paragraph{Codigo.}
Ver \texttt{cpp.tex}, seccion \textit{Closest pair of points}.

\vspace{0.5em}
